  \def\stat{med-gay}





\def\tit{ПРИМЕНЕНИЕ КЛАСТЕРИЗАЦИИ\\ В~ЗАДАЧАХ РАЗМЕЩЕНИЯ ПОДВИЖНЫХ ТОЧЕК 


ДОСТУПА\\ В~ВОЗДУШНО-НАЗЕМНЫХ БЕСПРОВОДНЫХ СЕТЯХ$^*$}





\def\titkol{Применение кластеризации в~задачах размещения подвижных точек 


доступа} % в~воздушно-наземных беспроводных сетях}








\def\aut{Е.\,Г.~Медведева$^1$, Э.\,М.~Хайров$^2$, Н.\,А.~Поляков$^3$, 


Ю.\,В.~Гайдамака$^4$}





\def\autkol{Е.\,Г.~Медведева, Э.\,М.~Хайров, Н.\,А.~Поляков, 


Ю.\,В.~Гайдамака}





\titel{\tit}{\aut}{\autkol}{\titkol}





\index{Е.\,Г.~Медведева$^1$, Э.\,М.~Хайров$^2$, Н.\,А.~Поляков$^3$, 


Ю.\,В.~Гайдамака$^4$}


\index{Stepchenkov Yu.\,A.}


\index{Morozov N.\,V.}


\index{Diachenko Yu.\,G.}


\index{Khilko D.\,V.}


\index{Stepchenko D.\,Yu.}








{\renewcommand{\thefootnote}{\fnsymbol{footnote}}


\footnotetext[1]


{Публикация подготовлена при поддержке Программы РУДН <<5-100>>. Исследование выполнено при 


финансовой поддержке РФФИ (проекты №\,18-07-00576 и~№\,20-07-01064).}} 








\renewcommand{\thefootnote}{\arabic{footnote}}


\footnotetext[1]{Российский университет дружбы народов; Институт проблем информатики Федерального исследовательского 


центра <<Информатика и~управление>> Российской академии наук, \mbox{medvedeva-eg@rudn.ru}}


\footnotetext[2]{Российский университет дружбы народов, emil.khayrov@gmail.com}


\footnotetext[3]{Российский университет дружбы народов, goto97@mail.ru}


\footnotetext[4]{Российский университет дружбы народов; Институт проблем информатики 


Федерального исследовательского центра <<Информатика и~управление>> Российской 


академии наук, \mbox{gaydamaka-yuv@rudn.ru}}





 \label{st\stat}





 \thispagestyle{headings}


 


 \vspace*{-12pt}  








     


  


      


  \Abst{Выполнен обзор задач, возникающих в~беспроводных сетях 


  с~подвижными базовыми станциями, размещенными на беспилотных 


летательных аппаратах (БПЛА). Сформулирована и~решена задача оптимизации для 


определения позиций БПЛА, обеспечивающих 


максимизацию вероятности покрытия в~области предоставления связи при 


ограничениях на интерференцию от соседних базовых станций. Для 


сравнительного анализа эффективности развертывания сети при применении 


адаптивной навигации выбраны два метода~--- метод k-сред\-них и~метод роя 


час\-тиц, в~основе которых лежит кластеризация пользователей. Применение 


методов проиллюстрировано для сценария обеспечения связью участников 


массового мероприятия на открытой местности (сценарий <<концерт>>).}


  


  \KW{БПЛА; воздушно-наземная сеть; метод роя частиц; метод k-сред\-них; 


вероятность покрытия}





\DOI{10.14357\!/\!08696527200403} 





\vskip 10pt plus 9pt minus 6pt





 \thispagestyle{headings}


 


 \vspace*{-12pt}


  


  \section{Введение}


  


  Для преодоления кризиса дефицита спектра частот, связанного с 


экспоненциальным ростом беспроводных устройств, одним из возможных 


решений представляется расширение спектра до более высоких диапазонов, 


в~которых по-прежнему доступно большое число лицензированных 


и~нелицензированных час\-тот. Многообещающим подходом стало 


повсеместное внедрение и~использование час\-тот миллиметровых волн 


(mmWave) в~диапазоне от~30 до~300~ГГц~[1]. С~этим несравненным 


преимуществом технология mmWave стала важной компонентой 


в~беспроводных системах~5G и~последующих сетях~5G (beyond 5G, B5G), 


харак-\linebreak\vspace*{-12pt}





\noindent


теризующихся требованиями к~высокой пропускной способности, 


широкой полосе пропускания, сверхбыстрым скоростям передачи данных, 


сверхнизкой задержке и~расширенным возможностям подключения для 


сетевого сообщества. В~будущих сетях~5G подразумевается поддержка 


большого разнообразия широкодиапазонных приложений, предусмотренных 


для множественных/многоканальных беспроводных систем (Multiple Gigabit 


Wireless Systems, MGWS), развернутых на нелицензированных частотах 


диапазона около 60 ГГц, зарезервированного для промышленных, научных и~медицинских целей, и~сетей WiGig (беспроводной гигабитной) как 


продолжение сетей Wi-Fi в~диапазоне 60~ГГц. Кроме того, использование 


миллиметрового диапазона длин волн выделено в~качестве ключевого решения 


сетей~5G <<Новое радио>> (New Radio, NR), разработанного консорциумом 


3GPP в~качестве унифицированного радиоинтерфейса для поддержки 


различных пользовательских услуг, таких как, например, усовершенствованная 


мобильная широкополосная связь (Enhanced Mobile Broadband, eMBB) 


и~сверхнадежная связь с низкой задержкой (Ultra-Reliable Low Latency 


Communication, URLLC). 


  


  Диапазон электромагнитных волн в~миллиметровом диапазоне, называемый 


также диапазоном сверхвысоких частот, определяется полосой спектра от~30 


до~300~ГГц, между микроволнами и~инфракрасными волнами. Наличие 


обильной незанятой полосы пропускания, доступной на частотах mmWave,~--- 


одно из ключевых преимуществ связи~5G~NR по сравнению с ограниченными 


ресурсами микроволнового спектра волн ниже~6~ГГц, которые в~настоящее 


время используются традиционными беспроводными системами 


и~существующими сетями~4G~LTE. Очевидно, что используемый 


микроволновый спектр волн ниже~6~ГГц для существующих беспроводных 


систем недостаточен для достижения гигабитных скоростей передачи данных. 


Значительное повышение скорости передачи данных было теоретически 


подтверждено теоремой Котельникова~[2], показывающей, что пропускная 


способность увеличивается линейно с ростом полосы пропускания. 


В~частности, доступный радиоресурс полос mmWave \mbox{71--76}  


и~81--86~ГГц (известный также, как E-band)~[3] значительно превосходит 


суммарно радиоресурс всех других полос лицензированного спектра. 


В~соответствии с~3GPP Release~16~[3] для раннего развертывания сетей~5G 


с~технологией mmWave были выбраны диапазоны частот 24,25--29,5  


и~37--43,5~ГГц. Кроме того, Международным союзом электросвязи (МСЭ) 


и~консорциумом 3GPP был разработан план исследований~[4] для стандартов~5G, 


охватывающих две полосы частот миллиметрового диапазона, которые будут 


использоваться для коммерческих нужд, а именно: 40 и~100~ГГц. 


Первоначально еще в~1979~г.\ для поддержки высокопроизводительной 


двухточечной связи МСЭ были выделены диапазоны частот E-band~---  


71--76 и~81--86~ГГц, призванные обеспечить полнодуплексное 


подключение со скоростью передачи данных~1~Гбит/с и~выше. 


 В~2015~г.\ Федеральная комиссия по связи США предложила 


гибкие правила обслуживания для полос частот mmWave, включая 28, 


37 и~39~ГГц для лицензированного использования и~64--71~ГГц для 


нелицензированного использования. Компании AT\&T и~Verizon для своих 


начальных развертываний сетей~5G сфокусированы на использовании полос 


частот~39 и~28~ГГц соответственно.


  


  \section{Обзор задач при~использовании подвижных точек доступа 


на~беспилотных летательных аппаратах}


  


  Перспективным решением для обеспечения пользователей беспроводных 


сетей широкополосным доступом в~местах отсутствия стационарных сотовых 


базовых станций стало приобретающее в~последние годы популярность 


применение передвижных базовых станций. Так, использование БПЛА, или дронов, и~размещения на них точек 


доступа (ТД) для предоставления услуг связи как в~лицензированном (LTE, 


WiFi), так и~в~нелицензируемом (mmWave) диапазоне частот обусловлено 


невероятной гибкостью и~скоростью развертывания такой сети. Беспилотные летательные аппараты, 


оснащенные внешними на\-прав\-лен\-ны\-ми антеннами с возможностью усиления 


сигнала, выступают не только как средства для инспекций, поиска и~фиксации 


чрезвычайных ситуаций при спасательных операциях в~удаленных местах, но 


и~являются частью сетевой транспортной инфраструктуры и~служат доступом 


к~услугам передачи данных для пользователей на поверхности земли при 


проведении всевозможных мероприятий на больших открытых  


площадках~[5--7]. 


  


  Особенность передачи данных на базе технологии mmWave с помощью 


\mbox{БПЛА} заключается в~необходимости наличия прямой видимости (Line-of-Sight, 


LoS) между ТД и~пользователем. Это жесткое требование компенсируется 


по\-движ\-ностью БПЛА, когда обеспечивается адаптивная навигация 


относительно пользователей, что позволяет значительно повысить 


эффективность передачи и~качество обслуживания~[8--11]. В~статьях 


рассмотрен ряд решений для эффективного размещения БПЛА, 


обеспечивающих оптимизацию услуг передачи данных конечным 


пользователям.


  


   Число исследований, посвященных применению БПЛА для поддержки 


беспроводных сетей~5G на базе технологии mmWave, превышает сотню. 


Благодаря возможности зависания над поверхностью земли, достаточной 


гибкости и~прос\-то\-те развертывания, высокой маневренности, быстрой 


реконфигурации сети, БПЛА могут быть размещены практически в~любых 


местах для выполнения функций воздушной ретрансляции на транзитных 


участках сети в~качестве базовой станции или в~качестве ТД. При этом, 


в~отличие от БПЛА с фиксированным крылом, у которых положение 


в~пространстве динамически изменяется во времени, положение беспилотных 


летательных винтокрылых аппаратов относительно фиксированно и~статично. 


В~этой связи исследование пространственной конфигурации сети БПЛА стало 


одной из важных проблем при проектировании беспроводных сетей. Для этого 


необходимо определять и~уметь оптимизировать пространственное положение и~траектории БПЛА в~соответствии с требованиями к производительности 


системы. Еще одной ключевой проблемой является возможность эффективного 


контроля развертывания нескольких БПЛА, или так называемая кластеризация 


БПЛА и~формирование роя БПЛА. Среди последних исследований 


пространственной конфигурации БПЛА с точки зрения оптимизации 


положения и~траектории известны работы~[12--14], вопросы управления 


развертыванием БПЛА освещены в~[15, 16].


  


  Основными преимуществами использования БПЛА с функцией сбора 


энергии и~возможностью хранения данных (кеширования) являются 


обеспечение непосредственной передачи кешированного популярного контента в~направлении наземных мобильных терминалов и~продление операционной 


работы БПЛА путем сбора энергии из окружающей среды (ветер, солнечная 


энергия). Первое свойство дополнительно позволяет снизить перегрузки на 


транзитных участках сети. Однако перерывы в~поступлении энергии в~силу 


непостоянности источника и~неопределенность кеширования создают 


дополнительные проблемы для надежного подключения и~повсеместного 


покрытия такими сетями. Вопрос координации и~взаимодействия между БПЛА и~наземными базовыми станциями описанного сценария рассмотрен  


в~\cite{15-md}. Кроме того, в~данной работе предложена кластеризации БПЛА, 


направленная на разгрузку трафика в~заданном ограниченном пространстве.


  


  Кеширование популярного контента на периферии сети (edge computing) 


в~периоды непиковой нагрузки представляется многообещающим решением 


как для уменьшения нагрузки транзитных соединений в~опорных сетях, так 


и~для повышения пропускной способности каналов и~улучшения качества 


восприятия услуг (Quality of Experience, QoE) пользователями. 


Мотивированные этими преимуществами, авторы~\cite{12-md} исследовали 


проблему развертывания БПЛА с~функцией кеширования, чтобы 


оптимизировать показатели качества восприятия пользователя QoE в~облачной 


сети радиодоступа (cloud radio access network, CRAN) путем упреждающей 


загрузки и~кеширования популярного контента в~непиковые часы. Для анализа 


авторами выбрана метрика QoE, характеризующаяся позицией <<человек 


в~контуре>>, которая фиксирует скорость передачи данных, задержку и~тип 


устройства каждого пользователя. В~этом случае сформулированная задача 


оптимизации заключается в~минимизации мощности передачи БПЛА при 


удовлетворении метрики QoE для каждого пользователя. Для решения этой 


проблемы предложен алгоритм прогнозирования с использованием концепции 


эхо-сетей (echo state network, ESN) для генерации функции распределения 


моментов времени запроса контента и~модели мобильности каждого наземного 


пользователя, а~для навигации БПЛА используется метод k-сред\-них.


  


  Для задач патрулирования, возникающих в~приложениях интернета вещей, 


например при сборе данных с датчиков контроля контура системы 


безопасности, применяются модели движения БПЛА, не зависящие от 


перемещения пользователей в~области предоставления связи. Модели 


патрулирования, как правило, представляют собой замкнутые траектории 


движения, цель которых состоит в~поддержании равномерного распределения 


воздушных базовых станций по рассматриваемой области. В~\cite{17-md} 


исследовано семейство траекторий, в~соответствии с которым мобильность 


БПЛА позволяет поддерживать однородный охват пользователей при заданной 


вероятности покрытия, а также получить преимущество с точки зрения 


энергоэффективности. В~част\-ности, показано, что для рассморенных частных 


случаев, а именно: радиальной спиралевидной навигации и~кольцевой овальной 


навигации,~--- среднее время замирания сигнала (average fade duration, AFD) 


уменьшается примерно на два порядка по сравнению со статическим случаем 


расположения БПЛА, например, в~соответствии с точечным пуассоновским 


пространственным поцессом. 


  


  Еще одно интересное приложение кластеризации дронов представлено в~\mbox{статье}~\cite{8-md}, 


  посвященной стратегиям развертывания БПЛА в~сельской 


мест\-ности. Авторы проводят исследование задачи оптимизации с учетом 


энерго\-эф\-фек\-тив\-ности и~без него и~сравнивают ско\-рость выполнения 


алгоритмов. Системная модель включает малоподвижных, почти статичных 


пользователей, размещенных внут\-ри прямоугольной области для трех типов 


пуассоновского точечного процесса~--- однородного, неоднородного 


и~кластерного. Для навигации \mbox{БПЛА} используются 4~алгоритма: 


геометрическая релаксация, метод \mbox{k-сред}\-них (k-means), метод k-сред\-них 


с~учетом энергоэффективности и~развертывание по сетке. Все алгоритмы 


максимизируют вероятность покрытия (\textit{англ}.\ coverage probability), 


представляющую собой долю пользователей в~области предостав\-ле\-ния связи, 


для которых соединение обеспечивают точки доступа на БПЛА. Ограничением 


решаемой алгоритмически задачи оптимизации размещения БПЛА служит 


запрет пересечения зон покрытия ТД, что обеспечивает снижение 


интерференции. В~работе проведено дополнительное исследование 


зависимости точности решения задачи от погрешности при определении 


координат пользователей.


{  %\looseness=1





} 


  


  \section{Примеры навигации беспилотных летательных аппаратов при~групповой мобильности 


пользователей}


  


  Решение упомянутых выше задач размещения ТД на БПЛА 


  с~по\-мощью методов кластеризации иллюстрируется ниже при сравнительном 


анализе двух методов определения позиции БПЛА~--- метода  


k-сред\-них~\cite{18-md} и~метода роя частиц (Particle Swarm Optimization, 


PSO)~\cite{10-md, 19-md}. Особенность состоит в~моделировании перемещения 


пользователей с~помощью модели групповой мобильности с~опорной точкой 


(Reference Point Group Mobility, RPGM). Модель предусматривает следование 


группы за условным лидером с одновременным случайным перемещением 


каждого отдельного объекта внутри группы в~ограниченной области на 


открытой местности~\cite{20-md} и~соответствующее динамическое изменение 


положений БПЛА для максимизации вероятности покрытия. Подключение 


пользователей к~широкополосной беспроводной сети организовано на базе 


технологии Wi-Fi (семейство протоколов IEEE~802.11, поддерживающих 


развертывание беспроводных локальных сетей Wireless Local Area Network, 


WLAN)~\cite{21-md}, в~которой в~качестве ТД выступают несколько 


\mbox{БПЛА} с~направленной антенной, обеспечивающих беспроводное 


радиопокрытие данной области. Физические параметры радиосреды 


и~приемопередающих устройств взяты из~\cite{10-md}. Показателем качества для 


оценки вероятности покрытия выбрано отношение сигнал/шум (signal-to-noise 


ratio, SNR), которое должно превосходить заданный стандартами 


порог~$\overline{\gamma}$.


  


  Исследован сценарий взаимодействия объектов беспроводной сети на БПЛА, 


соответствующий обеспечению связью пользователей при проведении 


мероприятия на открытой местности без крупных массивных сооружений 


(здания, деревья, опоры линий электропередачи), так называемый сценарий <<концерт>>~\cite{10-md}. На 


рис.~1 показано множество пользователей $\mathcal{U}$, число которых 


фиксировано и~равно $\vert \mathcal{U}\vert \hm=U$, и~$N_d$ устройств БПЛА 


на фиксированной высоте, при этом радиопокрытие каждого БПЛА 


определяется мощностью излучающей антенны и~характеризуется радиусом 


$R_d$, $d\hm=1, \ldots, N_d$~\cite{9-md}. 


  


  \begin{figure} %fig1


  \vspace*{1pt}


\begin{center}


\mbox{%


\epsfxsize=97.185mm  


\epsfbox{med-1.eps}


}


\end{center}


\vspace*{-9pt}


  \Caption{Сценарий взаимодействия пользователей сети БПЛА:


  \textit{1}~--- пользователь в~зоне действия БПЛА на высоте~$h_u$;


  \textit{2}~--- пользователь вне зоны действия БПЛА}


  \end{figure}


  


  Считаем, что пользователь $u\hm\in \mathcal{U}$ может установить по 


радиоканалу соединение c ТД на БПЛА~$d$, если $r_{(d,u)}\hm\leq R_d$, где 


$r_{(d,u)}$~--- расстояние между устройством пользователя и~центром зоны 


покрытия ТД на БПЛА. Вероятность покрытия вычисляется как доля 


пользователей в~области покрытия сети на БПЛА:





\vspace*{1pt}





\noindent


  \begin{equation*}


  p_c=\fr{1}{U}\sum\limits_{u\in \mathcal{U}} 1 \left\{ \gamma_{(d,u)}\geq 


\overline{\gamma}\right\}\,,\enskip d=1,\ldots , N_d\,,


 %\label{e1-md}


  \end{equation*}


  


  \vspace*{-3pt}


  


  \noindent


где $\gamma_{(d,u)}$~--- отношение сигнал\!/\!шум на устройстве пользователя 


$u$, принимающего сигнал от БПЛА~$d$, $u\hm\in \mathcal{U}$, $d\hm=1, 


\ldots , N_d$; $\overline{\gamma}$~--- заданный порог отношения сигнал/шум, 


определяющий качество принимаемого сигнала. 





  Позиции БПЛА в~области предоставления связи определяются как решение 


следующей задачи оптимизации:





\pagebreak





\noindent


  \begin{equation*}


  \sum\limits^{N_d}_{d=1} \sum\limits^U_{u=1} 1\left\{ r_{(d,u)}\leq 


R_d\right\}\to \max_{(x,y)}


%  \label{e2-md}


  \end{equation*}


при условии $r_{(i,u})+r_{(j,u)}\hm> R_i+R_j$, $i\not=j$, $j\hm=1, \ldots , N_d$.


%


Здесь $\mathbf{x}\hm=(x_1, \ldots , x_{N_d})$, $\mathbf{y}=(y_1,\ldots , y_{N_d})$; 


$(x_d,y_d)$~--- координаты центра зоны покрытия ТД на БПЛА~$d$, $d\hm=1, 


\ldots , N_d$.





  Согласно методу k-средних на первом шаге пользователи в~зависимости от 


местоположения разбиваются на кластеры по числу БПЛА. Следующий шаг~---


 построение диаграммы Вороного~\cite{22-md}, в~результате чего область 


предоставления связи разделяется на несколько подобластей, в~каждой из 


которых размещается БПЛА. Согласно методу роя частиц изменение позиции 


подвижной точки доступа подчиняется принципу наилучшего 


найденного в~пространстве положения БПЛА с точки зрения вероятности покрытия. 











  


  Исследование проведено для двух случаев разбиения пользователей на 


группы: (1)~группы одинакового размера; (2)~число пользователей в~группе 


имеет распределение Ципфа с параметром~0,5. Оба метода использованы для 


адаптивной


 навигации, предусматривающей об\-нов\-ле\-ние позиций БПЛА через


 фиксированные интервалы


 времени~$t_{\mathrm{upd}}$. В~пользу применения 


адаптивной навигации говорят результаты для усредненной по~300~с 


моделирования доле пользователей, находящихся в~зоне покрытия БПЛА, 


приведенные на рис.~2 для групп одина-\linebreak\vspace*{-12pt}





\begin{floatingfigure}{65mm} %fig2


\vspace*{1pt}


\begin{center}


\mbox{%


\epsfxsize=62mm  


\epsfbox{med-2.eps}


}


\end{center}


\vspace*{-9pt}


\Caption{Средняя доля пользователей в~зоне покрытия БПЛА: \textit{1}~--- пользователи в~зоне действия


одного БПЛА; \textit{2}~--- пользователи в~зоне действия более одного БПЛА; 


\textit{3}~--- число пользователей в~группе одинаковое; \textit{4}~--- число 


пользователей в~группе распределено по Ципфу}


\vspace*{-12pt}


\end{floatingfigure}





\noindent


кового размера. Здесь с точки зрения 


вероятности покрытия метод роя час\-тиц дает выигрыш до~60\%, а метод  


k-средн\-их~--- до~50\% по сравнению с~движением БПЛА без учета времени 


об\-нов\-ле\-ния~$t_{\mathrm{upd}}$, например при пат-\linebreak ру\-ли\-ро\-ва\-нии по заданным 


тра\-ек\-то\-риям. 


  








Из рис.~2 также видно, что метод роя частиц дает более высокую вероятность 


покрытия при сравнимых показателях интерференции. Последнее можно 


оценить по доле пользователей, находящихся в~зоне действия нескольких 


БПЛА и~испытывающих вследствие этого интерференцию, приводящую к 


снижению качества связи.





  Отдельной задачей является выбор длительности интервала~$t_{\mathrm{upd}}$. 


Рисунок~3,\,\textit{а} иллюстрирует снижение вероятности покрытия для 


случая, когда обновление координаты опорной точ-\hspace*{69.5mm}\linebreak\vspace*{-12pt}





\pagebreak





\vspace*{-12pt}





\





\begin{figure} %fig3


\vspace*{1pt}


\begin{center}


\mbox{%


\epsfxsize=128mm  


\epsfbox{med-3.eps}


}


\end{center}


\vspace*{-9pt}


\Caption{Вероятность покрытия~$p_c$ без учета~$t_{\mathrm{upd}}$~(\textit{a}) 


и~с~$t_{\mathrm{upd}}\hm=5$~с~(\textit{б}):


\textit{1}~--- PSO, равные группы; \textit{2}~--- PSO, неравные группы; \textit{3}~--- k-means, 


равные группы; \textit{4}~--- k-means, неравные группы}


\vspace*{-18pt}


\end{figure}











\begin{floatingfigure}{65mm} %fig4


  \vspace*{-1pt}


\begin{center}


\mbox{%


\epsfxsize=62mm  


\epsfbox{med-4.eps}


}


\end{center}


\vspace*{-9pt}


  \Caption{Вероятность покрытия при $t_{\mathrm{upd}}\hm=0{,}01$~с:


  \textit{1}~--- PSO; \textit{2}~--- k-means}


  \vspace*{9pt}


  \end{floatingfigure}


  


\noindent  


ки группы происходит 


слишком час\-то, что не позволяет БПЛА достичь целевой точки, полученной в~результате 


решения задачи оптимизации. На рис.~3,\,\textit{б} корректный 


выбор длины интервала $t_{\mathrm{upd}}$ ведет к росту вероятности покрытия, при этом 


множественные пики на графике соответствуют зна\-че\-ни\-ям 


вероятности покрытия непосредственно после достижения БПЛА новой 


позиции, полученной при применении соответствующего алгоритма.


 


  Зависимость интервала $t_{\mathrm{upd}}$ от относительной скорости перемещения 


пользователя и~БПЛА и~от длины волны исследована в~\cite{17-md}. 


В~рассмотренном случае для несущей частоты $f_c \hm=28$~ГГц длина волны 


$\lambda/2=5$~см, что при относительной ско\-рости до~3,6~м/с дает значение 


$t_{\mathrm{upd}}\hm=10$ мс. Однако, как видно из рис.~4, где приведены результаты 


моделирования для случая групп равного размера, выигрыш в~вероятности 


покрытия не превышает~4\%, что не сравнимо с существенно 


увеличивающимися затратами на пересчет координат БПЛА и~навигацию.


  


  \vspace*{-6pt}


  


  \section{Заключение}


  


  \vspace*{-4pt}


  


  Анализ статистических характеристик качества покрытия  


воз\-душ\-но-на\-зем\-ной беспроводной сетью области моделирования для 


сценария <<концерт>> с заданным порогом отношения сигнал\!/\!шум показал, 


что с точки зрения вероятности покрытия метод роя частиц для 


адаптивной навигации позволяет достичь более высоких значений вероятности 


покрытия, чем метод k-сред\-них. Интересной задачей в~дальнейшем видится 


исследование среднего времени замирания сигнала, представляющего 


собой среднее время, которое пользователь остается без связи из-за недостаточного качества радиоканала


 до ТД, обслуживающей 


соединение пользователя. Авторы планируют в~дальнейшем на базе 


разработанного симулятора расширить сценарии взаимодействия устройств, 


уделив внимание проблеме построения эффективной зоны покрытия 


в~воздушных беспроводных сетях на БПЛА с учетом расположения одной 


и~нескольких базовых станций, пространственного распределения 


пользователей внутри зоны покрытия сети и~управления распределением 


поступающей нагрузки от пользователей.


  


{\small\frenchspacing


{%\baselineskip=10.8pt


\addcontentsline{toc}{section}{Литература}


\begin{thebibliography}{99}


\bibitem{1-md}


\Au{Zhang L., Zhao H., Hou S., %Zhao~Z., Xu~H., Wu~X., 


\textit{et al.}} A~survey on 


5G millimeter wave communications for UAV-assisted wireless networks~/\!\!/ 


IEEE Access, 2019. Vol.~7. P.~117460--117504.


\bibitem{2-md}


\Au{Котельников В.\,А.} О пропускной способности эфира и~проволоки в~электросвязи: Мат-лы к I~Всесоюзному съезду по вопросам технической 


реконструкции дела связи и~развития слаботочной промышленности.~--- М.: 


Всесоюзный энергетический комитет, 1933. (Репринт статьи в~журнале 


УФН, 2006. Т.~176. Вып.~7. С.~762--770.)


\bibitem{3-md}


Technical Specifications and Technical Reports for a~5G based 3GPP system. 


Specification 21.916.


\bibitem{4-md}


5G Spectrum Recommendations, 5G Amer.~--- Bellevue, WA, USA, April 2017.


\bibitem{5-md}


\Au{Guillen-Perez A., Sanchez-Iborra~R., Cano~M.\,D., Sanchez-Aarnoutse~J.\,C., 


Garcia-Haro~J.} WiFi networks on drones~/\!\!/ 8th ITU Kaleidoscope Academic 


Conference: ICTs for a Sustainable World Proceedings.~--- 


Bangkok, Thailand, 2016. P.~1--8.


\bibitem{6-md}


\Au{Mozaffari M., Saad W., Bennis~M., Debbah~M.} Communications and control 


for wireless drone-based antenna array~/\!\!/ IEEE T. Commun., 2019. 


Vol.~67. Iss.~1. P.~820--834. doi: 10.1109/TCOMM.2018.2871453.


\bibitem{7-md}


\Au{Amer R., Saad W., Marchetti~N.} Mobility in the sky: Performance and 


mobility analysis for cellular-connected UAVs~/\!\!/ IEEE T. Commun., 


2020. Vol.~68. Iss.~5. P.~3229--3246. doi: 10.1109/TCOMM.2020.2973629.





\bibitem{10-md} %8


\Au{Hartigan A., Wong M.\,A.} Algorithm AS~136: A~k-means clustering 


algorithm~/\!\!/ J.~R. Stat. Soc.~C Appl., 1979. Vol.~28. Iss.~1. 


P.~100--108.


\bibitem{11-md} %9


\Au{Kennedy J., Eberhart R.} Particle swarm optimization~/\!\!/  IEEE  Conference  


(International) on Neural Networks IV Proceedings.~--- IEEE, 1995.  


P.~1942--1948.


\bibitem{9-md} %10


\Au{Zeng Y., Zhang R., Lim~T.\,J.} Wireless communications with unmanned aerial 


vehicles: Opportunities and challenges~/\!\!/ IEEE Commun. Mag., 2016. Vol.~54. 


Iss.~5. P.~36--42. 


\bibitem{8-md} %11


\Au{Sun J., Masouros C.} Deployment strategies of multiple aerial BSs for user 


coverage and power efficiency maximization~/\!\!/ IEEE T. Commun., 


2019. Vol.~67. Iss.~4. P.~2981--2994. doi: 10.1109/TCOMM.2018.2889460.





\bibitem{12-md} %12


\Au{Chen M., Mozaffari M., Saad W., Yin~C., Debbah~M., Hong~C.\,S.} Caching in 


the sky: Proactive deployment of cache-enabled unmanned aerial vehicles for 


optimized quality-of-experience~/\!\!/ IEEE J.~Sel. Area. Comm., 2017. Vol.~35. 


Iss.~5. P.~1046--1061. 


\bibitem{13-md}


\Au{Gapeyenko M., Bor-Yaliniz I., Andreev~S., Yanikomeroglu~H., 


Koucheryavy~Y.} Effects of blockage in deploying mmWave drone base stations for 


5G~Networks and beyond~/\!\!/ IEEE  Conference (International)  on 


Communications Workshops Proceedings.~--- IEEE, 2018. Art. ID: 


8403671. 6~p. doi: 10.1109/ICCW.2018.8403671.


\bibitem{14-md}


\Au{Ghazzai H., Ghorbel M.\,B., Kassler~A., Hossain~M.\,J.} Trajectory 


optimization for cooperative dual-band UAV swarms~/\!\!/ IEEE Global 


Communications Conference Proceedings.~--- IEEE, 2018. P.~1--7.


\bibitem{15-md}


\Au{Wu H., Tao X., Zhang N., Shen~X.} Cooperative UAV cluster-assisted 


terrestrial cellular networks for ubiquitous coverage~/\!\!/ IEEE J.~Sel. Area. 


Comm., 2018. Vol.~36. Iss.~9. P.~2045--2058. doi: 


10.1109/JSAC.2018.2864418.


\bibitem{16-md}


\Au{Khosravi Z., Gerasimenko~M., Andreev~S., Koucheryavy~Y.} Performance 


evaluation of UAV-assisted mmWave operation in mobility-enabled urban 


deployments~/\!\!/ 41st Conference (International) on Telecommunications 


and Signal Processing Proceedings.~--- IEEE, 2018. P.~150--153. doi: 


10.1109/TSP.2018.8441321.


\bibitem{17-md}


\Au{Enayati S., Saeedi H., Pishro-Nik~H., Yanikomeroglu~H.} Moving aerial base 


station networks: a stochastic geometry analysis and design perspective~/\!\!/ IEEE 


T. Wirel. Commun., 2019. Vol.~18. Iss.~6. P.~2977--2988.


\bibitem{18-md}


\Au{Tafintsev N., Gerasimenko M., Moltchanov~D., Akdeniz~M., Yeh~S., 


Himayat~N., Andreev~S., Koucheryavy~Y., Valkama~M.} Improved network 


coverage with adaptive navigation of mmWave-based drone-cells~/\!\!/ 


IEEE Global Communications Conference 


Proceedings.~---  IEEE, 2018. 


P.~1--7. doi: 


10.1109/GLOCOMW.2018.8644097.


\bibitem{19-md}


\Au{Xu D., Tian Y.} A~comprehensive survey of clustering algorithms~/\!\!/ Annals 


Data Science, 2015. Vol.~2. Iss.~2. P.~165--193. 


\bibitem{20-md}


\Au{Kalantari E., Bor-Yaliniz I., Yongacoglu~A., Yanikomeroglu~H.} User 


association and bandwidth allocation for terrestrial and aerial base stations with 


backhaul considerations~/\!\!/ 28th Annual Symposium (International) 


on Personal, Indoor, and Mobile Radio Communications Proceedings.~--- IEEE, 2017. P.~1--6. 


\bibitem{21-md}


\Au{Zolanvari M., Jain R., Salman~T.} Potential data link candidates for civilian 


unmanned aircraft systems: A~survey~/\!\!/ IEEE Commun. Surv. 


Tut., 2020. Vol.~22. Iss.~1. P.~292--319.


\bibitem{22-md}


\Au{Preparata F., Shamos M.} Computational geometry: An introduction.~--- 


Berlin--Heidelberg: Springer-Verlag, 1985. 390~p.


\end{thebibliography}


} }








\vspace*{-3pt}





\hfill{\small\textit{Поступила в~редакцию 12.09.20}}














%\vspace*{6pt}





%\hrule





%\vspace*{2pt}





\newpage





\vspace*{-28pt}





%\hrule





%\vspace*{8pt}





\def\tit{APPLICATION OF CLUSTERING IN~DEPLOYMENT OF~MOBILE ACCESS POINTS 


IN~AIR--GROUND WIRELESS NETWORKS}





\def\titkol{Application of clustering in~deployment of~mobile access points in


%~air--ground 


wireless networks}











\def\aut{E.\,G.~Medvedeva$^{1,2}$, E.\,M.~Khayrov$^1$, N.\,A.~Polyakov$^1$, 


and~Yu.\,V.~Gaidamaka$^{1,2}$}





\def\autkol{E.\,G.~Medvedeva, E.\,M.~Khayrov, N.\,A.~Polyakov, 


and~Yu.\,V.~Gaidamaka}





\titel{\tit}{\aut}{\autkol}{\titkol}





\vspace*{-11pt}








\noindent


$^1$Peoples' Friendship University of Russia (RUDN University), 6~Miklukho-\linebreak


$\hphantom{^1}$Maklaya Str.,  Moscow 117198, Russian Federation





\noindent


$^2$Institute of Informatics Problems, Federal Research Center ``Computer Science\linebreak


 $\hphantom{^1}$and Control''  


of the Russian Academy of Sciences, 44-2~Vavilov Str., Moscow\linebreak


$\hphantom{^1}$119333, Russian Federation





\def\leftfootline{\small{\textbf{\thepage}


\hfill Sistemy i Sredstva Informatiki~---


Systems and Means of Informatics 2020 vol~30 no\,4}


}%


 \def\rightfootline{\small{Sistemy i Sredstva Informatiki~---


 Systems and Means of Informatics 2020 vol~30 no\,4


\hfill \textbf{\thepage}}}


  


  


\vspace*{3pt}














\Abste{The paper provides an overview of tasks that arise in wireless networks with mobile base 


stations located on unmanned aerial vehicles (UAV). The authors choose two methods of adaptive 


navigation based on user clustering for a comparative analysis of effectiveness of network 


deployment~--- the k-means method and the particle swarm method. The solution to the 


optimization problem is the positions of UAV that maximize the probability of 


coverage in the communication provision area with restrictions on interference from neighboring 


base stations. The application of the methods is illustrated for the scenario of a~concert event which 


is defined as a~case for providing communication between participants of a~mass event in an open 


area.}





\KWE{UAV; air-ground network; aerial-terrestrial communication; particle swarm; k-means; 


coverage probability}











\DOI{10.14357\!/\!08696527200403} 





%\vspace*{-24pt}








\Ack


\noindent


The publication was prepared with the support of the ``RUDN University Program 5-100'' and 


funded by the RFBR according to the research projects No.\,18-07-00576 and No.\,20-07-01064. 


 





\vspace*{-12pt}





\renewcommand{\bibname}{\protect\rmfamily References}


%\renewcommand{\bibname}{\large\protect\rm References}





{\small\frenchspacing


{%\baselineskip=10.8pt


\addcontentsline{toc}{section}{References}


\begin{thebibliography}{99}


\bibitem{1-md-1}


\Aue{Zhang, L., H. Zhao, S.~Hou, %Z.~Zhao, H.~Xu, X.~Wu, 


\textit{et al.}} 2019. A survey on 5G 


millimeter wave communications for UAV-assisted wireless networks. \textit{IEEE Access} 


7:117460--117504.


\bibitem{2-md-1}


\Aue{Kotel'nikov, V.\,A.} 2006. On the transmission capacity of ether and wire in electric 


communications. \textit{Phys.-Usp.} 49(7):736--744.


\bibitem{3-md-1}


TS 21.916. 2018. Technical Specifications and Technical Reports for a 5G based 3GPP system. 


Available at: {\sf  


https://portal.3gpp.org/desktopmodules/Specifications/\linebreak 


SpecificationDetails.aspx?specificationId=3441} (accessed  


September~23, 2020).


\bibitem{4-md-1}


5G Spectrum Recommendations, 5~G Amer. April 2017. Bellevue, WA. Available at: {\sf 


 https://www.5gamericas.org/wp-content/uploads/2019/07/5GA\_5G\_\linebreak Spectrum\_Recommendations\_2017\_FINAL.pdf} 


(accessed  September~23, 2020).


\bibitem{5-md-1}


\Aue{Guillen-Perez, A., R.~Sanchez-Iborra, M.\,D.~Cano, J.\,C.~Sanchez-Aarnoutse, and 


J.~Garcia-Haro.} 2016. Wifi networks on drones. \textit{8th ITU Kaleidoscope Academic 


Conference: ICTs for a Sustainable World Proceedings}. Bangkok. 1--8.


\bibitem{6-md-1}


\Aue{Mozaffari, M., W. Saad, M.~Bennis, and M.~Debbah.} 2019. Communications and control 


for wireless drone-based antenna array. \textit{IEEE T. Commun.} 67(1):820--834. doi: 


10.1109/TCOMM.2018.2871453.


\bibitem{7-md-1}


\Aue{Amer, R., W. Saad, and N.~Marchetti.} 2020. Mobility in the sky: Performance and mobility 


analysis for cellular-connected UAVs. \textit{IEEE T. Commun.} 68(5):3229--3246. 


doi: 10.1109/TCOMM.2020.2973629.


\bibitem{10-md-1} %8


\Aue{Hartigan, A., and M.\,A.~Wong.} 1979. Algorithm AS 136: A~k-means clustering algorithm.  


\textit{J.~R. Stat. Soc.~C Appl.} 28(1):100--108.


\bibitem{11-md-1} %9


\Aue{Kennedy, J., and R. Eberhart.} 1995. Particle swarm optimization. \textit{IEEE  Conference 


(International) on Neural Networks IV Proceedings}. IEEE. 1942--1948.





\bibitem{9-md-1} %10


\Aue{Zeng, Y., R. Zhang, and T.\,J.~Lim.} 2016. Wireless communications with unmanned aerial 


vehicles: Opportunities and challenges. \textit{ IEEE Commun. Mag.} 54(5):36--42. 


\bibitem{8-md-1} %11


\Aue{Sun, J., and C. Masouros.} 2019. Deployment strategies of multiple aerial BSs for user 


coverage and power efficiency maximization. \textit{IEEE T.~Commun.} 


67(4):2981--2994. doi: 10.1109/TCOMM.2018.2889460.


\bibitem{12-md-1}


\Aue{Chen, M., M.~Mozaffari, W.~Saad, C.~Yin, M.~Debbah, and C.\,S.~Hong.} 2017. Caching 


in the sky: Proactive deployment of cache-enabled unmanned aerial vehicles for optimized 


 quality-of-experience. \textit{IEEE J.~Sel. Area. Comm.} 35(5):1046--1061. 


\bibitem{13-md-1}


\Aue{Gapeyenko, M., I.~Bor-Yaliniz, S.~Andreev, H.~Yanikomeroglu, and Y.~Koucheryavy.} 


2018. Effects of blockage in deploying mmWave drone base stations for 5G~networks and beyond. 


\textit{IEEE Conference (International) on Communications Workshops Proceedings}. IEEE. 


Art. ID:  8403671. 6~p.


doi: 10.1109/ICCW.2018.8403671.


\bibitem{14-md-1}


\Aue{Ghazzai, H., M.\,B.~Ghorbel, A.~Kassler, and M.\,J.~Hossain}. 2018. Trajectory 


optimization for cooperative dual-band UAV swarms. \textit{IEEE Global Communications Conference 


Proceedings}. IEEE. 1--7.


\bibitem{15-md-1}


\Aue{Wu, H., X.~Tao, N.~Zhang, and X.~Shen.} 2018. Cooperative UAV cluster-assisted 


terrestrial cellular networks for ubiquitous coverage. \textit{IEEE J.~Sel. Area. Comm.} 


36(9):2045--2058. doi: 10.1109/JSAC.2018.2864418.


\bibitem{16-md-1}


\Aue{Khosravi, Z., M. Gerasimenko, S.~Andreev, and Y.~Koucheryavy}. 2018. Performance 


evaluation of UAV-assisted mmWave operation in mobility-enabled urban deployments. 


\textit{41st Conference (International)  on Telecommunications and Signal Processing 


Proceedings}. IEEE. 150--153. doi: 


10.1109/TSP.2018.8441321.


\bibitem{17-md-1}


\Aue{Enayati, S., H. Saeedi, H.~Pishro-Nik, and H.~Yanikomeroglu.} 2019. Moving aerial base 


station networks: A~stochastic geometry analysis and design perspective. \textit{IEEE~T. 


Wirel. Commun.} 18(6): 2977--2988.


\bibitem{18-md-1}


\Aue{Tafintsev, N., M. Gerasimenko, D.~Moltchanov, M.~Akdeniz, S.~Yeh, N.~Himayat, 


S.~Andreev, Y.~Koucheryavy, and M.~Valkama.} 2018. Improved network coverage with adaptive 


navigation of mmWave-based drone-cells. \textit{IEEE Global Communications Conference Proceedings}. 


IEEE. 1--7.  doi: 


10.1109/GLOCOMW.2018.8644097.


\bibitem{19-md-1}


\Aue{Xu, D., and Y. Tian.} 2015. A~comprehensive survey of clustering algorithms. \textit{Annals 


Data Science} 2(2):165--193. 


\bibitem{20-md-1}


\Aue{Kalantari, E., I. Bor-Yaliniz, A.~Yongacoglu, and H.~Yanikomeroglu.} 2017. User 


association and bandwidth allocation for terrestrial and aerial base stations with backhaul 


considerations. \textit{28th Annual Symposium (International) on Personal, Indoor, and Mobile 


Radio Communications Proceedings}. IEEE. 1--6. 


\bibitem{21-md-1}


\Aue{Zolanvari, M., R. Jain, and T.~Salman.} 2020. Potential data link candidates for civilian 


unmanned aircraft systems: A~survey. \textit{IEEE Commun. Surv. Tut.} 


22(1):292--319.


\bibitem{22-md-1}


\Aue{Preparata, F., and M. Shamos.} 1985. \textit{Computational geometry: An introduction}. 


Berlin--Heidelberg: Springer-Verlag. 390~p.


\end{thebibliography}


} }





%\vspace*{-3pt}





\hfill{\small\textit{Received September 12, 2020}} 





%\vspace*{-16pt}





\Contr





\noindent


\textbf{Medvedeva Ekaterina G.} (b.\ 1985)~--- assistant professor, Department of Applied 


Probability and Informatics, Peoples' Friendship University of Russia (RUDN University), 


6~Miklukho-Maklaya Str., Moscow 117198, Russian Federation; senior scientist, Institute of 


Informatics Problems, Federal Research Center ``Computer Science and Control'' of the Russian 


Academy of Sciences, 44-2~Vavilov Str., Moscow 119333, Russian Federation; 


\mbox{medvedeva-eg@rudn.ru}





\vspace*{3pt}





\noindent


\textbf{Khayrov Emil M.} (b.\ 1997)~--- Master student, Department of Applied Probability and 


Informatics, Peoples' Friendship University of Russia (RUDN University), 6~Miklukho-Maklaya 


Str., Moscow 117198, Russian Federation; \mbox{emil.khayrov@gmail.com}








\vspace*{3pt}





\noindent


\textbf{Polyakov Nikita A.} (b.\ 1997)~--- Master student, Department of Applied Probability and 


Informatics, Peoples' Friendship University of Russia (RUDN University), 6~Miklukho-Maklaya 


Str., Moscow 117198, Russian Federation; \mbox{goto97@mail.ru}








\vspace*{3pt}





\noindent


\textbf{Gaidamaka Yuliya V.} (b.\ 1971)~--- Doctor of Science in physics and mathematics, 


professor, Department of Applied Probability and Informatics, Peoples' Friendship University of 


Russia (RUDN University), 6~Miklukho-Maklaya Str., Moscow 117198, Russian Federation; 


senior scientist, Institute of Informatics Problems, Federal Research Center ``Computer Science and 


Control'' of the Russian Academy of Sciences, 44-2~Vavilov Str., Moscow 119333, Russian 


Federation; \mbox{gaydamaka-yuv@rudn.ru}





\label{end\stat}





\renewcommand{\bibname}{\protect\rm Литература} 


